% Tabela de Metadados do Dataset Olímpico
\begin{table}[htbp]
\centering
\caption{Estrutura de metadados do dataset olímpico histórico}
\label{tab:dataset_metadata}
\footnotesize
\begin{tabular}{lp{2.5cm}p{8cm}}
\toprule
\textbf{Coluna} & \textbf{Tipo} & \textbf{Descrição} \\
\midrule
ID & Inteiro & Identificador único do atleta. Permite rastreamento longitudinal de participações em múltiplas edições olímpicas. \\
\midrule
Name & String & Nome completo do atleta conforme registrado nos arquivos olímpicos oficiais. \\
\midrule
Sex & Categórico (M/F) & Gênero do atleta: Masculino (M) ou Feminino (F). \\
\midrule
Age & Inteiro & Idade do atleta no momento da competição. Valores faltantes (NA) presentes para competidores históricos. \\
\midrule
Height & Numérico & Altura do atleta em centímetros. Valores faltantes (NA) comuns em registros anteriores a 1960. \\
\midrule
Weight & Numérico & Peso do atleta em quilogramas. Valores faltantes (NA) comuns em registros anteriores a 1960. \\
\midrule
Team & String & Nome completo da equipe/país representado pelo atleta (ex: "United States", "Brazil"). \\
\midrule
NOC & String (3 letras) & Código do Comitê Olímpico Nacional em formato padrão ISO (ex: "USA", "BRA", "CHN"). Campo mais confiável que Team para análises por país. \\
\midrule
Games & String & Identificador da edição olímpica no formato "Ano Temporada" (ex: "2016 Summer", "2014 Winter"). \\
\midrule
Year & Inteiro & Ano de realização dos jogos olímpicos (1896 a 2021). \\
\midrule
Season & Categórico & Temporada da competição: Verão (Summer) ou Inverno (Winter). \\
\midrule
City & String & Cidade sede dos jogos olímpicos (ex: "Athens", "Rio de Janeiro", "Beijing"). \\
\midrule
Sport & String & Modalidade esportiva (ex: "Swimming", "Basketball", "Football", "Athletics"). \\
\midrule
Event & String & Evento específico dentro da modalidade (ex: "Swimming Men's 100 metres Freestyle"). \\
\midrule
Medal & Categórico & Tipo de medalha conquistada: Gold, Silver, Bronze, ou NA para participações sem pódio. \\
\bottomrule
\end{tabular}
\fonte{\citeonline{Griffin2018OlympicDataset}}
\end{table}
